% =========================================================
% frontmatter/preface.tex
% Kata pengantar.
% =========================================================

\chapter*{Kata Pengantar}
\addcontentsline{toc}{chapter}{Kata Pengantar}
\markboth{Kata Pengantar}{Kata Pengantar}

Buku ajar ini disusun untuk mendukung mata kuliah
\textbf{\BookCourse} (\BookCourseCode) pada \BookProgram,
\BookFaculty, \BookInstitution. Mata kuliah ini berada pada ranah
pemrograman tingkat rendah, yaitu pemrograman yang berhadapan
langsung dengan arsitektur perangkat keras, register, memori, dan
instruksi mesin.

Pendekatan yang digunakan adalah \textit{Outcome-Based Education}
(OBE). Artinya, seluruh isi buku dirancang mundur dari capaian
pembelajaran: setiap bab menetapkan lebih dahulu kompetensi yang
harus dicapai mahasiswa, kemudian menyajikan materi, contoh,
praktikum, latihan, dan asesmen yang secara eksplisit mengukur
kompetensi tersebut. Dengan cara ini, materi perkuliahan, aktivitas
belajar, dan penilaian berada pada jalur yang sama dan dapat
dipertanggungjawabkan.

Bahasa assembly dipilih sebagai bahasa pengantar praktik karena
bahasa inilah yang paling dekat dengan instruksi mesin. Melalui
assembly, mahasiswa dapat melihat bagaimana sebuah algoritma
diterjemahkan menjadi urutan instruksi yang benar-benar dieksekusi
prosesor, bagaimana data berpindah antara register dan memori, serta
bagaimana kesalahan kecil pada level instruksi berdampak besar pada
perilaku program. Pemahaman ini menjadi fondasi yang kuat untuk
mata kuliah lanjutan seperti organisasi komputer, sistem operasi,
kompilator, dan keamanan perangkat lunak.

Buku ini disusun untuk dapat digunakan secara bertahap. Penulisannya
dilakukan per bab, dan setiap bab merupakan modul mandiri yang dapat
dipelajari serta dikompilasi secara terpisah. Rincian rencana
penulisan per bab dapat dilihat pada dokumen
\textit{Rencana Penulisan Buku Ajar} yang menyertai berkas buku ini.

Penulis menyadari bahwa buku ini masih memiliki keterbatasan.
Kritik, saran, dan koreksi dari pembaca sangat diharapkan untuk
perbaikan pada edisi berikutnya.

\vspace{1cm}

\begin{flushright}
\begin{tabular}{c}
Penulis\\[1.6cm]
\BookAuthor
\end{tabular}
\end{flushright}

\clearpage
